%%=============================================================================
%% Voorwoord
%%=============================================================================

\chapter*{\IfLanguageName{dutch}{Woord vooraf}{Preface}}%
\label{ch:voorwoord}

%% TODO:
%% Het voorwoord is het enige deel van de bachelorproef waar je vanuit je
%% eigen standpunt (``ik-vorm'') mag schrijven. Je kan hier bv. motiveren
%% waarom jij het onderwerp wil bespreken.
%% Vergeet ook niet te bedanken wie je geholpen/gesteund/... heeft

\chapter*{Voorwoord}
\label{ch:voorwoord}
\addcontentsline{toc}{chapter}{Voorwoord}

Voor u ligt de bachelorproef "Van abstracte data naar intuïtieve feedback: Het potentieel van Mixed Reality op de Meta Quest voor reanimatietraining binnen de 
geneeskunde". Dit onderzoek vormt het sluitstuk van mijn opleiding Professionele Bachelor in de Toegepaste Informatica aan HOGENT en is tot stand gekomen 
in het academiejaar 2025--2026.

Voor mij is technologie nooit een doel op zich geweest. 
Code, hardware en geavanceerde concepten zoals Mixed Reality zijn in mijn ogen puur een vehikel 
een middel om tastbare, reële problemen op te lossen en iets fundamenteels bij te dragen aan de maatschappij. 
Dit project lag mij bovendien nauw aan het hart; omdat ik in mijn directe omgeving van dichtbij heb meegemaakt welke impact cardiologische aandoeningen op een 
leven en een gezin kunnen hebben, kreeg dit onderzoek een diepere, persoonlijke betekenis. De ultieme motivator tijdens dit intensieve traject was dan ook de 
wetenschap dat deze applicatie een directe, positieve bijdrage kan leveren aan de kwaliteit van reanimatietraining. Iets bouwen wat in de toekomst kan helpen bij 
het effectiever aanleren van levensreddende handelingen: dat is wat mij heeft gedreven om dit reverse-engineering- en ontwikkeltraject tot een succes te brengen.

Dit project had echter niet gerealiseerd kunnen worden zonder de onmisbare steun en expertise van een aantal sleutelfiguren, die ik hier graag uitdrukkelijk wil bedanken.

Allereerst gaat mijn oprechte dank uit naar mijn promotor, Dhr. Jan Willem. Uw scherpe academische blik, constructieve kritiek en continue sturing hebben ervoor 
gezorgd dat dit onderzoek de juiste diepgang behield en dat de complexe technische materie helder en gestructureerd op papier werd gezet.

Daarnaast wil ik mijn co-promotor en opdrachtgever, Dhr. Thomas Cuelenaere, hartelijk bedanken. Uw diepgaande expertise binnen de gezondheidszorg en uw 
enthousiaste medewerking tijdens heel het traject waren van onschatbare waarde. U gaf mij niet alleen de ruimte om met dit innovatieve idee aan de slag te gaan, 
maar bood ook de onmisbare klinische inzichten om de softwarematige feedback perfect af te stemmen op de praktijk. 

Tot slot een bijzonder woord van dank aan mijn familie en vrienden. Jullie onvoorwaardelijke steun en het geloof in wat ik met dit project wilde bereiken, 
hebben me de rust gegeven om me volledig op dit onderzoek te kunnen focussen.

\vspace{1.5cm}

\noindent Dogukan Uyanik\\
Gent, mei 2026